% Condensed in rewrite 2026-09-25 into App. C.1 (short notes under tab:datasets); full original wording kept here.
\subsection{Datasets and splits}\label{app:data}
\Cref{tab:datasets} summarises the datasets, their splits and licences; the paragraphs below give the details.
\begin{table}[h]
\centering\scriptsize
\caption{\textbf{Datasets.} Every forecasting task predicts $K{=}10$ ($K{=}12$ on IWS) future feature grids from $H{=}3$ observed frames and
the actions. Step: one model time step. Splits are disjoint at the split unit. DINO-WM PushT and Wall use the official
releases and splits.}
\label{tab:datasets}
\setlength{\tabcolsep}{2.5pt}
\begin{tabular}{@{}p{0.14\linewidth}p{0.305\linewidth}ccccc@{}}
\toprule
Dataset & Recorded scenes & Episodes (train/val/test) & Step & Action & Split & Licence \\
\midrule
DROID (24 shards) & real Franka arm in homes, labs and offices; two exterior cameras & 851 / 143 / 132 & 0.33\,s & 35-D & session & CC-BY-4.0 \\
Open-H Hamlyn & dVRK arms on 7 surgical practice tasks (knot tying, needle grasp and handover, peg transfer, two suturing tasks, tissue lifting, tissue retraction); scene camera & 708 / 152 / 153 & 0.33\,s & 80-D & episode & CC-BY-4.0 \\
Language-Table & xArm pushing coloured blocks; fixed camera (Open-X subset) & 2{,}079 / 246 / 264 & 0.1\,s & 2-D & episode & CC-BY-4.0 \\
IWS PushT / Box / Rope & real bimanual T-block, box and rope tasks \citep{wang2026iws} & 510--512 / 90 / 200 & 0.17\,s & 20/70/40-D & recording & as released \\
PushT / TwoRoom / Reacher & simulated T-block pushing, two-room navigation and two-link reaching (LeWM planning suites) & 2{,}607 / 290 (PushT) & 5 env steps & 10-D & episode & MIT \\
\bottomrule
\end{tabular}
\end{table}

\paragraph{DROID.} We use 24 official full-release shards (22.0\,GB) containing 1{,}126 episodes from 14 sites (327{,}498 native steps, 64{,}826 action blocks; 316/59/58 recording
sessions in train/val/test). Splits are assigned by recording session with a salted hash, so no session appears in two
splits. Frames are subsampled every 5 native steps (15\,Hz $\to$ 3\,Hz), and each action block
concatenates the 5 commanded Cartesian poses and gripper positions ($5\times7{=}35$).
Models are trained on the first exterior camera; the second exterior camera records the same episodes at the same
instants with the same action blocks from a different viewpoint; the main models never see it in training and are
evaluated on it zero-shot.
\paragraph{Open-H-Embodiment Hamlyn.} Seven dVRK training tasks, released under CC-BY-4.0 in LeRobot
format. We use the \texttt{observation.images.color} stream, which the release labels ``endoscope''
but which is an Intel RealSense scene camera ($848{\times}480$). One step is 1/3\,s: stride 10 at
30\,Hz, or stride 5 for the 15\,Hz tissue-lifting task. Each action block holds five commands at
15\,Hz, each an 8-D pose-and-gripper target for each arm ($5\times16{=}80$). Episodes shorter than 14
steps (6 of 1{,}019) are removed. Splits are 70/15/15 per task, assigned by a hash of the episode id.
\paragraph{Language-Table.} The real-robot Language-Table release \citep{lynch2023interactive} in Open-X RLDS form:
an xArm with a cylindrical end effector pushes rigid, coloured blocks on a wooden board, seen by a fixed camera
($640{\times}360$, resized to 224 like Bridge and RT-1 and shown at its native aspect). We fetch 14 of the 1{,}024 pinned
shards (6.0\,GB, 6{,}161 episodes). The Open-X sheet lists 10\,Hz control, so one model step is one native step (0.1\,s)
and the action is the 2-D delta end-effector setpoint. We do not subsample further: the episodes are short
hindsight-labelled segments of longer continuous play (median 12 native steps), so a 0.3\,s step would leave only 331 episodes long enough for
$H{+}K{=}13$ steps, and the end effector already moves $\approx$2.7\,cm per native step (median). Episodes shorter than 14
steps are dropped (3{,}572), leaving 2{,}589, split 10\%/9\%/81\% into test/val/train by a salted hash of the episode id.

Open-X ships only a
\texttt{train} split, so every id keeps that name (\texttt{train00160\_052}: shard 160, record 52) and we carve
validation and test from it. The split is by segment: the release has no id of the source play episode, so segments
of one play session can fall into different splits and share a scene layout. In our subset we find no direct
continuations: for only 1 of 264 test segments does a training segment start or end with a frame closer to the test
segment's boundary frame than the median one-step frame change,
but source-level disjointness is not guaranteed.
\paragraph{Planning suites.} PushT, TwoRoom and Reacher expert datasets from the LeWM release
(HDF5). Episodes are split disjointly into train and validation sets for world-model training.
Evaluation follows the official protocol (\cref{app:planning}).
